% ====== 第六部分 外观设计国际申请 ======
\cleardoublepage\setguidefootmark{6}{第六部分}\resetchapters
\thispagestyle{plainempty}\hypertarget{gpos.p6}{}\begingroup\centering\guidecenterpage
\dadati\fontsize{29.04bp}{43bp}\selectfont
\vspace*{242bp}
第六部分\par
外观设计国际申请\par
\phantomsection
\addcontentsline{toc}{part}{第六部分\quad 外观设计国际申请}
\endgroup\clearpage
\cleardoublepage\pagestyle{tocpage}
\guideoutline[1]{分目录}{parttoc6}
\gdef\guidetabcontent{\hbox{第}\hbox{六}\hbox{部}\hbox{分}}
{\centering\guidecenterpage\heiti\fontsize{14.04bp}{21bp}\selectfont 目\quad 录\par}
\vspace{6bp}
\printparttoc{p6}
\clearpage\pagestyle{main}
\startparttoc{p6}
\guidechapter{gpos.6.1}{第一章}{外观设计国际注册申请的事务处理}{}{0bp}
\guidesection{gpos.6.1.1}{1.}{引言}{}{74.4bp}
本章涉及申请人按照工业品外观设计国际注册海牙协定日内瓦文本（以下简称海牙协定）通过专利局提交外观设计国际注册申请、缴纳费用，以及按照海牙协定提出并指定中国的外观设计国际注册申请（以下简称外观设计国际申请）的手续审查和事务处理的特别规定。本章没有说明和规定的事项，参照本指南第一部分、第五部分的规定。

\guidesection{gpos.6.1.2}{2.}{外观设计国际注册申请的提交}{}{102.0bp}
\guidesubsection{gpos.6.1.2.1}{2.1}{提交途径}{}{0bp}
外观设计国际注册申请可以直接向国际局提交。申请人在中国有经常居所或者营业所的，可以通过专利局向国际局提交外观设计国际注册申请。

通过专利局提交外观设计国际注册申请的，国际程序中后续其他文件应当直接向国际局提交。

\guidesubsection{gpos.6.1.2.2}{2.2}{传送与不传送}{}{102.0bp}
\guidesubsubsection{gpos.6.1.2.2.1}{2.2.1}{收到日的确定}{}{0bp}
通过专利局向国际局提交的外观设计国际注册申请，如果国际局于专利局收到日起1个月内收到，以专利局收到日视为国际局收到日，否则以国际局实际收到之日为收到日。

\guidesubsubsection{gpos.6.1.2.2.2}{2.2.2}{传送条件}{}{74.4bp}
外观设计国际注册申请符合下列条件的，专利局予以传送国际局：

{\setlength{\parindent}{17.1bp}（1）至少有申请人之一在中国有经常居所或者营业所；\par}

{\setlength{\parindent}{17.1bp}（2）至少有申请人之一选择中国作为申请人的缔约方；\par}

{\setlength{\parindent}{17.1bp}（3）使用英语撰写外观设计国际注册申请文件；\par}

{\setlength{\parindent}{17.1bp}（4）使用海牙协定规定的正式表格；\par}

{\setlength{\parindent}{17.1bp}（5）申请中包括外观设计图片或者照片；\par}

{\setlength{\parindent}{17.1bp}（6）包含中国内地中文通信信息；\par}

{\setlength{\parindent}{17.1bp}（7）申请文件中不得包含违反法律、社会公德或者妨害公共利益的信息。\par}

外观设计国际注册申请指定中国的，申请人可以提交外观设计国际注册申请的中文译文。

\guidesubsection{gpos.6.1.2.3}{2.3}{传送与不传送程序}{}{102.0bp}
\guidesubsubsection{gpos.6.1.2.3.1}{2.3.1}{文件处理}{}{0bp}
专利局收到外观设计国际注册申请后，进行如下文件处理：

{\setlength{\parindent}{17.1bp}（1）确定收到日：申请人通过专利局提交外观设计国际注册申请的，以专利局实际收到日为收到日；\par}

{\setlength{\parindent}{17.1bp}（2）给出提交编号：专利局按照收到外观设计国际注册申请的先后顺序给出提交编号。\par}

\guidesubsubsection{gpos.6.1.2.3.2}{2.3.2}{传送程序}{}{74.4bp}
外观设计国际注册申请符合传送条件的，传送程序如下：

{\setlength{\parindent}{17.1bp}（1）向申请人发出外观设计国际注册申请传送通知书，告知传送编号、传送期限及文件清单；\par}

{\setlength{\parindent}{17.1bp}（2）向国际局传送外观设计国际注册申请的文件及收到日等数据。\par}

\guidesubsubsection{gpos.6.1.2.3.3}{2.3.3}{不传送程序}{}{74.4bp}
外观设计国际注册申请不符合传送条件的，向申请人发出外观设计国际注册申请不传送通知书，告知申请人不传送的原因。

向专利局当面递交的外观设计国际注册申请，不符合传送条件的，应当直接向当事人说明原因，不予接收。

\guidesection{gpos.6.1.3}{3.}{外观设计国际申请的事务处理}{}{102.0bp}
\guidesubsection{gpos.6.1.3.1}{3.1}{确定在中国的申请日}{}{0bp}
按照海牙协定已确定国际注册日并指定中国的外观设计国际申请，视为向专利局提出的外观设计专利申请，该国际注册日视为专利法第二十八条所称的申请日。

\guidesubsection{gpos.6.1.3.2}{3.2}{给予国家申请号}{}{74.4bp}
外观设计国际申请经国际局公布后，专利局对国际局传送的外观设计国际申请给予国家申请号并进行后续审查。

\guidesubsection{gpos.6.1.3.3}{3.3}{其他文件的受理}{}{102.0bp}
\guidesubsubsection{gpos.6.1.3.3.1}{3.3.1}{其他文件的受理条件}{}{0bp}
外观设计国际申请经国际局公布后，外观设计国际申请的当事人向专利局办理相关手续时，应当使用中文提交符合规定的相关文件，写明国家申请号，并根据专利法第十八条的规定办理委托手续。

其他规定适用本指南第五部分第三章第3.1节的规定。

\guidesubsubsection{gpos.6.1.3.3.2}{3.3.2}{其他文件的受理程序}{}{74.4bp}
适用本指南第五部分第三章第3.2节的规定。

\guidesubsection{gpos.6.1.3.4}{3.4}{分案申请受理}{}{74.4bp}
针对外观设计国际申请提出的分案申请，除符合本指南第五部分第三章第2.3.2.1节的规定外，还应当核实分案申请请求书中是否填写了原申请的申请日和原申请的申请号，该原申请的申请日应当是其国际注册日，原申请的申请号填写原申请的国际注册号。该分案申请按照国家申请处理。

\guidesubsection{gpos.6.1.3.5}{3.5}{公告程序}{}{74.4bp}
对外观设计国际申请作出给予保护的决定后，专利局予以公告，公告的内容包括：专利权的著录事项以及一幅图片或者照片。著录事项主要包括：分类号、专利号、国际注册号、授权公告号（出版号）、申请日、授权公告日、优先权事项、专利权人事项、使用该外观设计的产品名称等。公告著录事项内容在国际注册公布文本中已有记载的，与其保持一致。该外观设计专利权自公告之日起在中国生效。专利局公告后，外观设计国际申请的申请人可以请求专利局出具外观设计国际申请专利登记簿副本，作为在中国给予保护的证明。外观设计专利的单行本的内容包括扉页、图片或者照片以及简要说明。其中，图片或者照片、简要说明以国际局公布的给予保护声明所确定的文本形式提供。

国际局已公告的其他事项，除涉及权利转移的以外，以国际局公告为准。

\guidesubsection{gpos.6.1.3.6}{3.6}{相关手续审查}{}{102.0bp}
\guidesubsubsection{gpos.6.1.3.6.1}{3.6.1}{著录项目变更}{}{0bp}
外观设计国际申请的申请人（或者专利权人）权利变更、名称和/或地址变更、在国际局的代理事项变更的，当事人应当向国际局办理相关手续。

外观设计国际申请的申请人（或者专利权人）权利变更的，当事人除了向国际局办理相关手续外，还应当按照专利法实施细则的规定向专利局提交证明文件，证明文件适用本指南第一部分第一章第6.7.2.2节和第6.7.2.6节的规定，证明文件是外文的，应当同时附具中文题录译文。没有提交证明文件或者提交证明文件不合格的，专利局应当通知国际局该权利变更在中国未生效。

\guidesubsubsection{gpos.6.1.3.6.2}{3.6.2}{权利的恢复}{}{74.4bp}
外观设计国际申请的当事人，因未及时答复驳回通知，其专利申请被视为撤回的，当事人可按照本指南第五部分第七章第6节的相关规定请求恢复权利。

\guidesubsubsection{gpos.6.1.3.6.3}{3.6.3}{未续展的专利权终止}{}{74.4bp}
外观设计国际申请在专利局授权公告之后，专利权人未按照海牙协定的规定办理续展手续的，专利权自在中国的申请日起满5年或者10年之日起终止。

\guidesubsubsection{gpos.6.1.3.6.4}{3.6.4}{权利的部分放弃}{}{74.4bp}
外观设计国际申请在专利局授权公告之后，专利权人向国际局提出针对中国放弃部分权利的，该部分放弃的生效日为国际局登记之日。

\guidesection{gpos.6.1.4}{4.}{缴费的特别规定}{}{102.0bp}
\guidesubsection{gpos.6.1.4.1}{4.1}{国际程序费用的缴纳}{}{0bp}
外观设计国际注册申请的国际程序相关费用应当直接向国际局缴纳。通过专利局提交外观设计国际注册申请文件的，可以通过专利局向国际局缴纳外观设计国际注册申请的相关费用。

通过专利局缴纳费用的，当事人应当以传送编号为依据，通过网上缴费或者直接向专利局当面缴纳相关费用。缴纳费用时应注明正确的传送编号以及缴纳的费用名称。不符合上述规定的，视为未办理缴费手续。

通过专利局缴纳的外观设计国际注册申请相关费用全部转交国际局，国际局以其账户收到费用的日期为缴费日。国际局收取外观设计国际申请的单独指定费后转交专利局。专利局不进行上述相关费用的退费。国际程序中费用相关事宜，由当事人直接与国际局联系。

\guidesubsection{gpos.6.1.4.2}{4.2}{专利局费用的缴纳}{}{74.4bp}
外观设计国际申请经国际局公布之后，当事人向专利局缴纳相关费用的，应当以国家申请号或者国际注册号缴费。

\guidechapter{gpos.6.2}{第二章}{外观设计国际申请的审查}{}{0bp}
\guidesection{gpos.6.2.1}{1.}{引言}{\guidenote{细则 143}}{74.4bp}
外观设计国际申请的审查是指专利局根据专利法及其实施细则的规定对申请人按照海牙协定提出并指定中国的外观设计国际注册申请的审查。根据专利法实施细则第一百四十三条的规定，外观设计国际申请经专利局审查后没有发现驳回理由的，由专利局作出给予保护的决定，并通知国际局。

\guidenote{细则 50.1}本章涉及的专利局对外观设计国际申请的审查范围：

{\setlength{\parindent}{17.1bp}（1）明显实质性缺陷的审查，包括外观设计国际申请是否明显属于专利法第五条第一款、第二十五条第一款第（六）项规定的情形，是否不符合专利法实施细则第十一条的规定，是否明显不符合专利法第二条第四款、第二十三条第一款和第二款、第二十七条第二款、第三十一条第二款、第三十三条、专利法实施细则第一百四十二条的规定，或者依照专利法第九条的规定不能取得专利权。\par}

{\setlength{\parindent}{17.1bp}（2）其他文件和相关手续的审查，包括与外观设计国际申请有关的其他文件和相关手续是否符合专利法第十八条、第二十四条，以及专利法实施细则第三条第一款、第十八条、第三十三条第四款、第三十四条第三款、第一百三十九条、第一百四十条、第一百四十一条的规定。\par}

本章仅对上述审查中的特殊问题作出说明和规定。与国家申请相同的其他问题，本章没有说明和规定的，参照本指南第一部分第三章、第四部分第五章、第五部分第十章的规定。

\guidesection{gpos.6.2.2}{2.}{审查原则}{}{74.4bp}
对于外观设计国际申请，应当依据以下原则进行审查：

{\setlength{\parindent}{17.1bp}（1）申请的形式或者内容适用海牙协定以及1999年文本和1960年文本共同实施细则的规定，审查员不得以申请文件的形式缺陷为由驳回外观设计国际申请。\par}

{\setlength{\parindent}{17.1bp}（2）明显实质性缺陷以及其他文件和相关手续的审查，适用专利法及其实施细则和本指南的规定。\par}

\guidesection{gpos.6.2.3}{3.}{审查程序}{}{102.0bp}
\guidesubsection{gpos.6.2.3.1}{3.1}{作出给予保护的决定}{\guidenote{细则 138 及}}{0bp}
\guidenote{143.1}外观设计国际申请经审查没有发现驳回理由的，审查员应当作出给予保护的决定，向国际局发出给予保护声明。给予保护的外观设计国际申请包括不需要向国际局发出驳回通知就符合授权条件的国际申请，以及答复驳回通知后符合授权条件的国际申请。

\guidesubsection{gpos.6.2.3.2}{3.2}{作出驳回通知}{}{74.4bp}
外观设计国际申请存在明显实质性缺陷的，审查员应当向国际局发出驳回通知。

驳回通知应当包含驳回所依据的全部理由以及相应的法律条款。如果驳回理由涉及专利法第二十三条第一款、第二款或者专利法第九条的规定的，还应当包含与该外观设计国际申请相关的现有设计或者国内同样的外观设计专利申请或者专利的相关信息。

\guidesubsection{gpos.6.2.3.3}{3.3}{驳回通知的答复}{\guidenote{细则 3.1}}{74.4bp}
申请人在收到驳回通知后，应当在指定的期限内根据专利法第十八条的规定办理委托手续，并进行答复。根据专利法实施细则第三条第一款的规定，申请人进行答复时应当使用中文提交陈述意见，使用英文对申请文件进行修改。

对于答复文件中出现新的缺陷的，如果该缺陷可以通过补正克服，审查员应当进行全面审查，并向申请人发出补正通知书；如果该缺陷是不能通过补正方式克服的明显实质性缺陷，审查员应当向申请人发出审查意见通知书。

\guidesubsection{gpos.6.2.3.4}{3.4}{作出驳回决定}{\guidenote{细则 50.2}}{74.4bp}
申请人针对驳回通知或者审查意见通知书提交的答复文件未能克服通知书中指出的明显实质性缺陷的，审查员可以作出驳回决定。

{\setlength{\parindent}{22.2bp}驳回决定内容适用本指南第一部分第三章第3.5节的规定。\par}

\guidesubsection{gpos.6.2.3.5}{3.5}{前置审查与复审后的处理}{}{74.4bp}
适用本指南第一部分第三章第3.6节的规定。

\guidesection{gpos.6.2.4}{4.}{审查依据的文本确认}{}{102.0bp}
\guidesubsection{gpos.6.2.4.1}{4.1}{审查依据的文本}{}{0bp}
作为外观设计国际申请审查基础的文本可能包括：

{\setlength{\parindent}{17.1bp}（1）国际局公布的外观设计国际申请的英文文本。\par}

{\setlength{\parindent}{17.1bp}（2）根据专利法实施细则第一百四十一条提交的修改文本。\par}

{\setlength{\parindent}{17.1bp}（3）根据专利法实施细则第五十条提交的英文补正文本。\par}

\guidesubsection{gpos.6.2.4.2}{4.2}{国际局公布的外观设计国际申请文件的效力}{}{74.4bp}
根据海牙协定第14条第（1）款的规定，外观设计国际申请自国际注册日起享有与在中国提出的外观设计专利申请同等的效力。

\guidesection{gpos.6.2.5}{5.}{外观设计国际申请文件的审查}{}{102.0bp}
\guidesubsection{gpos.6.2.5.1}{5.1}{著录项目的审查}{}{0bp}
外观设计国际申请的著录项目以国际局公布的为准，审查员一般不对其进行审查，但申请人为克服专利局发出的通知书中指出的缺陷而修改了著录项目的情形除外。

\guidesubsection{gpos.6.2.5.2}{5.2}{图片或者照片的审查}{}{102.0bp}
\guidesubsubsection{gpos.6.2.5.2.1}{5.2.1}{视图名称及其标注}{}{0bp}
外观设计国际申请的视图名称及其标注被视为符合本指南第一部分第三章第4.2.1节的规定。

\guidesubsubsection{gpos.6.2.5.2.2}{5.2.2}{图片或者照片的清楚表达}{\guidenote{法 27.2}}{74.4bp}
专利法第二十七条第二款规定，申请人提交的有关图片或者照片应当清楚地显示要求专利保护的产品的外观设计。

审查员应当就申请人提交的有关图片或者照片是否存在影响要求保护的产品整体或者局部外观设计清楚表达的明显实质性缺陷进行审查。

\guidesubsection{gpos.6.2.5.3}{5.3}{简要说明书的审查}{\guidenote{细则 142}}{74.4bp}
根据专利法实施细则第一百四十二条的规定，国际局公布的外观设计国际申请中包括含设计要点的说明书的，视为已经依照专利法实施细则第三十一条的规定提交了简要说明。

\guidenote{法 64.2}根据专利法第六十四条第二款的规定，外观设计国际申请简要说明书的内容用于解释图片或者照片所示的产品的外观设计。审查员应当结合简要说明书的内容和产品名称，就外观设计图片或者照片是否清楚地表达了要求保护的产品整体或者局部的外观设计进行审查。

\guidesubsection{gpos.6.2.5.4}{5.4}{根据专利法第五条第一款、第二十五条第一款第（六）项、专利法实施细则第十一条和专利法第二条第四款的审查}{}{74.4bp}
适用本指南第一部分第三章第6节、第7节的规定。

\guidesubsection{gpos.6.2.5.5}{5.5}{根据专利法第九条和第二十三条第一款、第二款的审查}{}{74.4bp}
适用本指南第一部分第三章第11节、第8节的规定。

\guidesubsection{gpos.6.2.5.6}{5.6}{根据专利法第三十一条第二款的审查}{}{74.4bp}
对于外观设计国际申请，审查员应当审查其是否符合专利法第三十一条第二款的规定。

一件外观设计国际申请包括两项以上（含两项）外观设计的，申请人可以主动提出或者依据审查员的审查意见提出分案申请。分案申请被视为国家申请。

\guidenote{细则 141}根据专利法实施细则第一百四十一条的规定，申请人主动提出分案申请的，应当自外观设计国际申请公布之日起两个月内提出。

申请人按照审查员的审查意见提出分案申请的，最迟应当在原申请的国内公告日起两个月内提出。上述期限届满后，或者原申请已被驳回，或者原申请被视为撤回且未被恢复权利的，一般不得再提出分案申请。

涉及分案的其他规定适用本指南第一部分第三章第9.4节的规定。

\guidesubsection{gpos.6.2.5.7}{5.7}{根据专利法第三十三条的审查}{}{74.4bp}
适用本指南第一部分第三章第10.2节的规定。

\guidesection{gpos.6.2.6}{6.}{其他文件和相关手续的审查}{}{102.0bp}
\guidesubsection{gpos.6.2.6.1}{6.1}{委托专利代理机构}{}{0bp}
外观设计国际申请的申请人答复驳回通知或者办理其他专利事务时，应当符合专利法第十八条第一款、专利法实施细则第十八条的有关规定。

如果申请人在提交外观设计国际申请时已经委托了符合专利法第十八条规定的专利代理机构，在向专利局办理专利事务时，需要按照本指南第一部分第一章第6.1.2节的规定办理委托手续。

解除和辞去委托的规定适用本指南第一部分第一章第6.1.3节的规定。

\guidesubsection{gpos.6.2.6.2}{6.2}{优先权的审查}{}{74.4bp}
根据专利法实施细则第一百三十七条的规定，外观设计国际申请的国际注册日视为专利法第二十八条所称的申请日。

除本节特殊规定外，优先权的其他规定参照本指南第一部分第三章第5.2节的规定。外观设计国际申请已提出优先权要求并被国际局接受的，不收取优先权要求费。

\guidesubsubsection{gpos.6.2.6.2.1}{6.2.1}{要求外国优先权}{}{102.0bp}
\guideparagraph{gpos.6.2.6.2.1.1}{6.2.1.1}{在先申请和要求优先权的在后申请}{}{0bp}
申请人要求以其在先提出的外观设计国际申请为基础享有优先权的，视为根据专利法第二十九条第一款的规定要求外国优先权。

\guideparagraph{gpos.6.2.6.2.1.2}{6.2.1.2}{要求优先权声明}{\guidenote{细则 139}}{74.4bp}
根据专利法实施细则第一百三十九条的规定，国际局公布的外观设计国际申请中包括一项或者多项优先权的，视为已经依照专利法第三十条的规定提出了书面声明。

\guideparagraph{gpos.6.2.6.2.1.3}{6.2.1.3}{在先申请文件副本}{\guidenote{细则 139}}{74.4bp}
根据专利法实施细则第一百三十九条的规定，外观设计国际申请的申请人要求优先权的，应当自外观设计国际申请公布之日起三个月内向专利局提交第一次提出的专利申请的副本。在先申请文件副本中可以不包含该副本的中文题录译文。期满未提交的，视为未要求优先权。

\guideparagraph{gpos.6.2.6.2.1.4}{6.2.1.4}{在后申请的申请人}{\guidenote{细则 34.3}}{74.4bp}
根据专利法实施细则第三十四条第三款的规定，在后申请的申请人与在先申请文件副本中记载的申请人不一致的，申请人应当在外观设计国际申请公布之日起三个月内向专利局提交相关的证明文件。申请人期满未提交的，视为未要求优先权。

\guidesubsubsection{gpos.6.2.6.2.2}{6.2.2}{要求本国优先权}{}{102.0bp}
\guideparagraph{gpos.6.2.6.2.2.1}{6.2.2.1}{在先申请和要求优先权的在后申请}{}{0bp}
申请人要求以其在先在中国提出的外观设计为基础享有优先权的，视为根据专利法第二十九条第二款的规定要求本国优先权。

在后外观设计国际申请的国际注册日之前，专利局已经对在先申请发出授予专利权通知书和办理登记手续通知书，并且申请人已经办理了登记手续的，在后申请视为未要求优先权。

\guideparagraph{gpos.6.2.6.2.2.2}{6.2.2.2}{要求优先权声明}{\guidenote{细则 139}}{74.4bp}
根据专利法实施细则第一百三十九条的规定，国际局公布的外观设计国际申请中包括一项或者多项优先权的，视为已经依照专利法第三十条的规定提出了书面声明。

\guideparagraph{gpos.6.2.6.2.2.3}{6.2.2.3}{在先申请文件副本}{}{74.4bp}
适用本指南第一部分第三章第5.2.2.3节的规定。

\guideparagraph{gpos.6.2.6.2.2.4}{6.2.2.4}{在后申请的申请人}{\guidenote{细则 34.3}}{74.4bp}
根据专利法实施细则第三十四条第三款的规定，在后申请的申请人与在先申请文件副本中记载的申请人不一致的，申请人应当在外观设计国际申请公布之日起三个月内向专利局提交相关的证明文件。申请人期满未提交的，视为未要求优先权。

\guideparagraph{gpos.6.2.6.2.2.5}{6.2.2.5}{视为撤回在先申请的程序}{\guidenote{细则 35.3}}{74.4bp}
{\setlength{\parindent}{22.2bp}要求一项或者多项本国优先权的在后外观设计国际申请，经初步审查认为符合规定的，如果相应的在先申请是外观设计专利申请，且未办理登记手续，则该在先外观设计申请被视为撤回。\par}

被视为撤回的在先申请不得请求恢复。

\guidesubsubsection{gpos.6.2.6.2.3}{6.2.3}{优先权要求的撤回}{}{74.4bp}
申请人不得向专利局提出撤回优先权要求。

\guidesubsubsection{gpos.6.2.6.2.4}{6.2.4}{优先权要求的恢复}{}{74.4bp}
外观设计国际申请被视为未要求优先权的，不予恢复。

\guidesubsection{gpos.6.2.6.3}{6.3}{不丧失新颖性的公开}{\guidenote{细则 140}}{74.4bp}
根据专利法实施细则第一百四十条的规定，外观设计国际申请涉及的外观设计有专利法第二十四条第（二）项或者第（三）项所列情形的，应当在提出外观设计国际申请时声明，并在自外观设计国际申请公布之日起两个月内向专利局提交有关符合规定的证明文件。

\guidenote{细则 33.4}根据专利法实施细则第三十三条第四款的规定，外观设计国际申请涉及的外观设计有专利法第二十四条第（一）项或者第（四）项所列情形的，专利局认为必要时，可以要求申请人在指定期限内向专利局提交证明文件。

审查员应当对证明材料中注明的相关日期以及内容与请求保护的外观设计是否明显相关进行审查。

\stopparttoc
